% commentaire.tex — f174r, vv. 959–1000 (conversion initiale des travaux antérieurs)
\comm{982}{διπλῆν.}{Pollux (IV, 105) compte la διπλῆ parmi les figures de danse pour la tragédie (τραγικῆς ὀρχήσεως σχήματα) :\citComm{Καὶ μὲν τραγικῆς ὀρχήσεως σχήματα σιμή χείρ, καλαθίσκος, χεὶρ καταπρανής, ξύλου παράληψις, \textbf{διπλῆ}, θερμαυστρίς, κυβίστησις, παραβῆναι τὰ τέτταρα}Hésychius (δ 1948) ne nous apprend rien davantage :\citComm{διπλῆ· ὀρχήσεως εἶδος, ἢ κρούματος.}Il nous est difficile de reconstituer, même grossièrement, la chorégraphie de ce choral puisqu'Aristophane recourt au vocabulaire technique de la danse de son époque, dont on ne sait plus grand chose.}
